% Where the optimum of a convex problem sits: five contour plots over one polyhedron.
%
%   figures/tikz/ci-contours-convex.tex
%       -> media/figures/ci-contours-convex.{png,svg}
%
% Lecture 18 (Introduction to Constrained Optimization), "A taxonomy of contour plots".
% Drawn 2026-10-07 for the figures Biegler shows as Figs. 4.1 and 4.2 (pp. 66-67), which the
% lecture had only described in words. Our own geometry, not a redrawing of his.
%
% GEOMETRY (1 unit = 1 cm before scaling). One polyhedron F in every panel, the one of
% ra-polyhedron.tex: x1 >= 0, x2 >= 0, x1 + 2 x2 <= 4, 3 x1 + x2 <= 6, with vertices
% (0,0), (2,0), (1.6,1.2), (0,2).
%   (1) LP, f = -(x1 + 1.2 x2). Vertex values of x1 + 1.2 x2: 0, 2, 3.04, 2.4, so the
%       unique minimizer is the vertex (1.6,1.2). Contours x1 + 1.2 x2 = 1, 2, 3.04, 4.
%   (2) LP, f = -(3 x1 + x2). Vertex values 0, 6, 6, 2: the whole edge (2,0)-(1.6,1.2)
%       is optimal (it lies on 3 x1 + x2 = 6). Contours 3 x1 + x2 = 2, 4, 6, 7.5.
%   (3)-(5) f = |x - c|^2 (convex quadratic; circles as contours, so the minimizer over F
%       is the Euclidean projection of c onto F and can be checked by hand).
%   (3) c = (2.2,1.7). c - (1.6,1.2) = (0.6,0.5) = 0.18 (1,2) + 0.14 (3,1), a nonnegative
%       combination of the two outward normals at that vertex, so the projection is the
%       vertex (1.6,1.2); touching radius |(0.6,0.5)| = 0.781.
%   (4) c = (1.8,0.6) + 0.6 (0.949,0.316) = (2.369,0.790): the edge midpoint plus 0.6 times
%       the unit outward normal of 3 x1 + x2 <= 6, so the projection is the midpoint, in the
%       interior of that edge; touching radius 0.6.
%   (5) c = (0.8,0.8), inside F (distances to the four lines 0.8, 0.8, 0.716, 0.885), so the
%       minimizer is c itself and the constraints are inactive.
% In every panel the contour through the minimizer is solid; the others are dashed.
%
% GREYSCALE: F is hatched; contours are grey and dashed except the touching one (solid);
% the minimizer is a filled black dot, an optimal edge a thick black segment, and "+" marks
% the unconstrained minimizer c. Nothing is carried by colour alone.
\usetikzlibrary{patterns}
\definecolor{okabeblue}{HTML}{0072B2}

\begin{tikzpicture}[
    scale=0.80,
    >={Latex[length=1.6mm]},
    font=\small,
    lab/.style={font=\scriptsize, inner sep=1pt},
    cont/.style={black!55, dashed, line width=0.5pt},
    touch/.style={black, line width=0.7pt},
  ]
  \def\cifeas{%
    \path[pattern=north east lines, pattern color=okabeblue!55]
       (0,0) -- (2,0) -- (1.6,1.2) -- (0,2) -- cycle;
    \draw[okabeblue, line width=1.2pt] (0,0) -- (2,0) -- (1.6,1.2) -- (0,2) -- cycle;
    \node[lab, fill=white] at (0.35,1.5) {$\mathcal{F}$};}
  \def\ciframe{\draw[black!30, line width=0.4pt] (-0.4,-0.4) rectangle (2.9,2.6);}
  % ---------------------------------------------------------------- (1) LP, vertex
  \begin{scope}[xshift=0cm]
    \ciframe
    \begin{scope}
      \clip (-0.4,-0.4) rectangle (2.9,2.6);
      \cifeas
      \foreach \c in {1.0,2.0,4.0}
        \draw[cont] (-0.6,{(\c+0.6)/1.2}) -- (3.1,{(\c-3.1)/1.2});
      \draw[touch] (-0.6,{(3.04+0.6)/1.2}) -- (3.1,{(3.04-3.1)/1.2});
    \end{scope}
    \fill (1.6,1.2) circle (2.4pt);
    \draw[->, line width=0.7pt] (2.3,0.5) -- ++(0.32,0.38);
    \node[lab, anchor=north] at (2.42,0.47) {$-\nabla f$};
    \node at (1.25,-0.85) {(1)};
  \end{scope}
  % ---------------------------------------------------------------- (2) LP, edge
  \begin{scope}[xshift=3.7cm]
    \ciframe
    \begin{scope}
      \clip (-0.4,-0.4) rectangle (2.9,2.6);
      \cifeas
      \foreach \c in {2.0,4.0,7.5}
        \draw[cont] ({(\c+0.6)/3},-0.6) -- ({(\c-2.8)/3},2.8);
      \draw[touch] ({(6+0.6)/3},-0.6) -- ({(6-2.8)/3},2.8);
    \end{scope}
    \draw[line width=2.4pt] (2,0) -- (1.6,1.2);
    \draw[->, line width=0.7pt] (2.25,1.75) -- ++(0.57,0.19);
    \node[lab, anchor=south] at (2.5,1.95) {$-\nabla f$};
    \node at (1.25,-0.85) {(2)};
  \end{scope}
  % ---------------------------------------------------------------- (3) quadratic, vertex
  \begin{scope}[xshift=7.4cm]
    \ciframe
    \begin{scope}
      \clip (-0.4,-0.4) rectangle (2.9,2.6);
      \cifeas
      \draw[cont] (2.2,1.7) circle (0.35);
      \draw[touch] (2.2,1.7) circle (0.781);
      \draw[cont] (2.2,1.7) circle (1.2);
    \end{scope}
    \node[font=\normalsize] at (2.2,1.7) {$+$};
    \fill (1.6,1.2) circle (2.4pt);
    \node at (1.25,-0.85) {(3)};
  \end{scope}
  % ---------------------------------------------------------------- (4) quadratic, edge interior
  \begin{scope}[xshift=11.1cm]
    \ciframe
    \begin{scope}
      \clip (-0.4,-0.4) rectangle (2.9,2.6);
      \cifeas
      \draw[cont] (2.369,0.790) circle (0.3);
      \draw[touch] (2.369,0.790) circle (0.6);
      \draw[cont] (2.369,0.790) circle (1.0);
    \end{scope}
    \node[font=\normalsize] at (2.369,0.790) {$+$};
    \fill (1.8,0.6) circle (2.4pt);
    \node at (1.25,-0.85) {(4)};
  \end{scope}
  % ---------------------------------------------------------------- (5) quadratic, interior
  \begin{scope}[xshift=14.8cm]
    \ciframe
    \begin{scope}
      \clip (-0.4,-0.4) rectangle (2.9,2.6);
      \cifeas
      \draw[cont] (0.8,0.8) circle (0.3);
      \draw[cont] (0.8,0.8) circle (0.6);
      \draw[cont] (0.8,0.8) circle (1.0);
    \end{scope}
    \fill (0.8,0.8) circle (2.4pt);
    \node at (1.25,-0.85) {(5)};
  \end{scope}
  % group headings
  \node[font=\small] at (3.1,3.0) {linear $f$};
  \draw[black!50] (-0.4,2.75) -- (6.6,2.75);
  \node[font=\small] at (12.35,3.0) {convex quadratic $f$};
  \draw[black!50] (7.0,2.75) -- (17.7,2.75);
\end{tikzpicture}
